%! Author = 50669
%! Date = 3/18/2026

% Preamble
\documentclass[11pt]{article}

% Packages
\usepackage{amsmath}

% Document
\begin{document}
    \section*{Question 1}
    First let
    \[
        x_j=-1+jh,\qquad h=\frac{2}{n+1},\qquad j=0,1,\dots,n+1,
    \]
    and let $u\in C^4([-1,1])$ with $u(-1)=u(1)=0$, defining
    \[
        V_j:=\frac{u(x_{j-1})-2u(x_j)+u(x_{j+1})}{h^2}, \qquad j=1,\dots,n
    \]
    Then using Taylor expansion around xj,
    \begin{align*}
        u(x_{j+1}) &= u(x_j)+h u'(x_j)+\frac{h^2}{2}u''(x_j)+\frac{h^3}{6}u^{(3)}(x_j)+\frac{h^4}{24}u^{(4)}(\xi_j^+),\\
        u(x_{j-1}) &= u(x_j)-h u'(x_j)+\frac{h^2}{2}u''(x_j)-\frac{h^3}{6}u^{(3)}(x_j)+\frac{h^4}{24}u^{(4)}(\xi_j^-),
    \end{align*}
    for some $\xi_j^\pm$ between $x_{j-1}$ and $x_{j+1}$.\\
    Then adding both lines and subtracting $2u(x_j)$
    \[
        u(x_{j-1})-2u(x_j)+u(x_{j+1})=h^2u''(x_j)+\frac{h^4}{24}\big(u^{(4)}(\xi_j^-)+u^{(4)}(\xi_j^+)\big).
    \]
    Hence
    \[
        V_j-u''(x_j)=\frac{h^2}{24}\big(u^{(4)}(\xi_j^-)+u^{(4)}(\xi_j^+)\big).
    \]
    Since $u^{(4)}$ is continuous, so there exist $M>0$ such that
    $|u^{(4)}(x)|\le M$ for all $x\in[-1,1]$. Therefore
    \[
        |V_j-u''(x_j)|\le \frac{h^2}{24}(M+M)=\frac{M}{12}h^2.
    \]
    Let $K=\frac{M}{12}$, then take the maximum over interior indices:
    \[
        \max_{1\le j\le n}|V_j-u''(x_j)|\le Kh^2.
    \]
    So the differentiation matrix is second-order accurate in $\|\cdot\|_\infty$.

    \section*{Question 7:}
    First Crank-Nicholson scheme:
    \[
        u_{i+1} = u_i + \frac{\tau}{2}\left[f(t_{i+1}, u_{i+1}) + f(t_i, u_i)\right]
    \]
    then substitute $f(t,y) = Ay + g(t)$, that
    \[
        u_{i+1} = u_i + \frac{\tau}{2}\left[Au_{i+1} + g(t_{i+1}) + Au_i + g(t_i)\right]
    \]
    then rearrange the terms, put $u_{i+1}$ on left and $u_i$ on the right
    \begin{align*}
        u_{i+1} = u_i + \frac{\tau}{2}Au_{i+1} + \frac{\tau}{2}g(t_{i+1}) + \frac{\tau}{2}Au_i + \frac{\tau}{2}g(t_i) \\
        => \qquad u_{i+1} - \frac{\tau}{2}Au_{i+1} = u_i + \frac{\tau}{2}Au_i + \frac{\tau}{2}g(t_{i+1})  + \frac{\tau}{2}g(t_i) \\
        => \qquad (I-\frac{\tau}{2}A) u_{i+1} = (I+\frac{\tau}{2}A) u_i + \frac{\tau}{2}g(t_{i+1})  + \frac{\tau}{2}g(t_i)
    \end{align*}
    so we can extract:
    \begin{align}
        B = I - \frac{\tau}{2}A \\
        L = I + \frac{\tau}{2}A \\
        f_{i} = \frac{\tau}{2}[g{t_{i+1}} + g(t_{i})]
    \end{align}
    Thus, with the initial condition is $u_0 = y_0$, the Crank-Niocholson schmem can be written as:
    \[
        Bu_{i+1} = Lu_{i} + f_{i}
    \]


\end{document}